% LAFS: Load-Aware Flow Scheduler
% COMP-6910: Services, Semantic, and Cloud Computing
% Memorial University of Newfoundland — Group 7 — Winter 2026
% 10-page LNCS format

\documentclass[runningheads]{llncs}

\usepackage{graphicx}
\usepackage{amsmath}
\usepackage{amssymb}
\usepackage{booktabs}
\usepackage{hyperref}
\usepackage{xcolor}
\usepackage{listings}
\usepackage{microtype}
\usepackage{caption}
\usepackage{subcaption}
\usepackage{algorithm}
\usepackage{algpseudocode}
\usepackage{multirow}
\usepackage{tabularx}

\hypersetup{
  colorlinks=true,
  linkcolor=blue,
  citecolor=blue,
  urlcolor=blue
}

\lstset{
  basicstyle=\ttfamily\small,
  breaklines=true,
  frame=single,
  numbers=left,
  numberstyle=\tiny,
  keywordstyle=\color{blue},
  commentstyle=\color{gray}
}

\begin{document}

\title{LAFS: Load-Aware Flow Scheduling in Data Centre Fat-Tree Networks Using MILP and Hybrid Load Prediction}

\author{Victor Chisom Muojeke\inst{1} \and
Chiemerie Cletus Obijiaku\inst{1} \and
Olaleye Adeniyi Akinuli\inst{1}}

\authorrunning{Muojeke et al.}

\institute{Memorial University of Newfoundland, St. John's NL A1C 5S7, Canada\\
\email{\{v\_muojek, c\_obijai, o\_akinul\}@mun.ca}}

\maketitle

\begin{abstract}
We present LAFS (Load-Aware Flow Scheduler), a hybrid centralised-distributed scheduling system for data-centre fat-tree networks. LAFS combines a per-link hybrid load predictor (EWMA + ARIMA) with a Mixed-Integer Linear Program (MILP) that minimises maximum link utilisation across a scheduling window. Evaluated on a simulated $k{=}8$ fat-tree (128 hosts, 80 switches) with 1\,000 Facebook web-search flows at four load levels (30–90\%), LAFS achieves MILP-optimal placement in 100\% of scheduling instances (average solve time 3.0\,s, PuLP/CBC), matching ECMP flow completion times while providing a formal guarantee on maximum link utilisation ($z^* = 0.241$ at 50\% load). The full source code, 486 passing unit and integration tests, and the comparison experiment runner are publicly available.
\end{abstract}

\keywords{Flow scheduling \and Fat-tree networks \and MILP \and Load prediction \and ECMP \and Data centre networking}

%% =========================================================
\section{Introduction}
\label{sec:intro}
%% =========================================================

Modern data-centre fabrics carry heterogeneous traffic mixes: short ``mice'' flows from web-serving and RPC chains that are latency-sensitive, and long ``elephant'' flows from MapReduce shuffles and ML training that dominate link bandwidth. Equal-Cost Multi-Path routing (ECMP) hashes flows uniformly across available paths, producing well-known hash collisions that concentrate elephant flows onto shared links and inflate tail flow-completion times (FCT)~\cite{hedera,conga}.

Two complementary strategies have emerged. \emph{Centralised schedulers} (Hedera~\cite{hedera}) detect elephant flows and re-route them via a global-knowledge controller. \emph{Distributed schedulers} (CONGA~\cite{conga}) measure per-path congestion locally and steer flowlets without global coordination. Neither approach provides a formal optimality guarantee on the resulting link-utilisation profile.

We propose LAFS, which integrates both concerns: a lightweight prediction module supplies short-horizon link-load forecasts, and a MILP formulation produces an assignment of flows to candidate ECMP paths that minimises the worst-case link utilisation with a 1\% optimality gap. The key contributions are:

\begin{enumerate}
  \item A hybrid EWMA+ARIMA load predictor with inverse-error blending and automatic $\alpha$ tuning that produces point estimates with 90\% confidence intervals.
  \item A compact MILP formulation that scales to $k{=}8$ fat-trees with 1\,000 concurrent flows using only open-source solvers (PuLP/CBC).
  \item A comprehensive empirical comparison of ECMP, Hedera, CONGA, and LAFS on a realistic Facebook web-search workload with hot-spot traffic.
  \item A fully reproducible implementation: 486 passing tests, a CLI experiment runner, and all result artefacts committed to a public repository.
\end{enumerate}

%% =========================================================
\section{Background and Related Work}
\label{sec:background}
%% =========================================================

\subsection{Fat-Tree Topology}

Al-Fares et al.~\cite{fattree} introduced the $k$-ary fat-tree, a three-tier Clos network built from commodity $k$-port switches. A $k{=}8$ fat-tree hosts $k^3/4 = 128$ servers across $5k^2/4 = 80$ switches. Each pod contains $k/2 = 4$ edge and $k/2 = 4$ aggregation switches; $k^2/4 = 16$ core switches interconnect pods. The topology provides $k/2 = 4$ edge-disjoint paths between any two hosts in different pods, offering full bisection bandwidth when traffic is distributed evenly.

\subsection{Baseline Schedulers}

\textbf{ECMP} assigns each flow to a path determined by a 5-tuple CRC32 hash. Its $O(1)$ per-flow cost makes it the universal baseline; its weakness is hash collisions that co-locate elephants.

\textbf{Hedera}~\cite{hedera} polls the network periodically, identifies elephant flows ($\geq$\,100\,MB in our parameterisation), and re-assigns them using Global First Fit (GFF): each elephant is placed on the least-loaded path. Mice flows remain on ECMP paths.

\textbf{CONGA}~\cite{conga} implements distributed, flowlet-level load balancing. Each switch maintains a Distributed Resilience table (DRE) updated by congestion feedback in packet headers. LAFS uses a software simulation of the DRE mechanism with a 100\,$\mu$s flowlet gap and exponential DRE decay ($\tau = 0.2$\,s).

\subsection{Relevant Recent Work}

The Alibaba HPN paper~\cite{alibabahpn} describes hyperscale RDMA fabrics for LLM training, where flow completion time variance directly impacts GPU utilisation. Their architecture motivates the need for scheduler mechanisms that provide latency guarantees beyond what ECMP hash-based routing can offer, lending relevance to MILP-guaranteed placement.

%% =========================================================
\section{System Model and Assumptions}
\label{sec:model}
%% =========================================================

\subsection{Network Model}

We model the data-centre network as a directed graph $G = (V, E)$ where $V$ is the set of switches and hosts, and each directed edge $(u,v) \in E$ has capacity $c_{uv}$. Host-edge links have capacity $C_h = 1$\,Gbps; fabric links (edge--aggregation, aggregation--core) have capacity $C_f = 10$\,Gbps, reflecting a standard 10:1 oversubscription at the host level.

\subsection{Flow Model}

A flow $f$ is characterised by its source host $s_f$, destination host $d_f$, and byte size $B_f$. We adopt the Facebook web-search distribution from Benson et al.~\cite{benson}: 80\% of flows are mice ($<$\,1\,MB), carrying 20\% of total bytes; the remaining 20\% are elephants. Flows arrive according to a Poisson process with rate $\lambda$ adjusted to achieve a target load fraction $\rho \in \{0.30, 0.50, 0.70, 0.90\}$ of the average fabric link capacity.

\subsection{Scheduling Window}

The scheduler operates in batches of duration $T_s = 0.1$\,s. Within each window it has access to the set of active flows and the most recent link-load forecast. For evaluation we simulate $T_{\text{sim}} = 10$\,s with $n = 1{,}000$ flows and measure steady-state metrics.

\subsection{Performance Metrics}

\begin{itemize}
  \item \textbf{FCT} — flow completion time; we report P50, P95, P99 percentiles.
  \item \textbf{Slowdown} — $\text{FCT} / \text{FCT}_{\text{alone}}$, normalising for flow size.
  \item \textbf{Fabric max utilisation} — $\max_{(u,v) \in E_f} \rho_{uv}$ where $E_f \subset E$ excludes host-edge links.
  \item \textbf{Jain's Fairness Index}~\cite{jain} — $\mathcal{F} = (\sum_f \text{tput}_f)^2 / (n \sum_f \text{tput}_f^2)$.
\end{itemize}

%% =========================================================
\section{Method: LAFS Design}
\label{sec:method}
%% =========================================================

LAFS has three components: a load predictor, a MILP flow placer, and a scheduling loop. Figure~\ref{fig:architecture} shows the high-level architecture.

\subsection{Hybrid Load Predictor}

For each fabric link $l$, LAFS maintains a time series of per-window utilisation measurements $\{u_l^{(t)}\}$. Two models are fitted online:

\medskip
\noindent\textbf{EWMA Predictor.} The exponentially-weighted moving-average estimate is
\[
  \hat{u}_l^{(t+1)}_{\text{EWMA}} = \alpha\, u_l^{(t)} + (1-\alpha)\, \hat{u}_l^{(t)}_{\text{EWMA}},
\]
where $\alpha$ is tuned online by minimising mean-squared one-step-ahead error over a sliding window of 20 samples.

\medskip
\noindent\textbf{AR Predictor.} A $p$-th order autoregressive model
\[
  \hat{u}_l^{(t+1)}_{\text{AR}} = \sum_{i=1}^{p} \phi_i\, u_l^{(t+1-i)} + \epsilon,
\]
fitted by ordinary least squares via QR decomposition. The order $p \in \{1,\ldots,5\}$ is selected by AIC. If the \texttt{statsmodels} ARIMA fitter is available, a full ARIMA(1,1,1) model is used instead.

\medskip
\noindent\textbf{Hybrid Blending.} The final forecast uses inverse-error weighting:
\[
  w_{\text{EWMA}} = \frac{1/e_{\text{EWMA}}}{1/e_{\text{EWMA}} + 1/e_{\text{AR}}}, \quad
  \hat{u}_l^{(t+1)} = w_{\text{EWMA}}\,\hat{u}_{l,\text{EWMA}} + (1-w_{\text{EWMA}})\,\hat{u}_{l,\text{AR}},
\]
where $e_{\text{EWMA}}, e_{\text{AR}}$ are exponentially-smoothed mean-absolute errors. A 90\% confidence interval is derived by bootstrapping the residual distribution.

\subsection{MILP Flow Placement}

Let $\mathcal{F}$ be the set of flows in the current scheduling window and $\mathcal{P}_f$ the set of candidate ECMP paths for flow $f$ (at most 16 in a $k{=}8$ fat-tree). Let $x_{f,p} \in \{0,1\}$ indicate whether flow $f$ is assigned to path $p$, and let $z \in [0,1]$ be the maximum link utilisation.

\medskip
\noindent\textbf{Objective.}
\begin{equation}
  \min_{x,\,z}\;\; z \;+\; \lambda \sum_{f \in \mathcal{F}_m} \sum_{p \in \mathcal{P}_f} (h_p - 1)\,x_{f,p}
  \label{eq:milp-obj}
\end{equation}
where $\mathcal{F}_m \subset \mathcal{F}$ is the set of mice flows, $h_p$ is the hop count of path $p$, and $\lambda = 10^{-3}$ is a small penalty that biases mice flows toward shorter paths without materially affecting $z$.

\medskip
\noindent\textbf{Assignment constraints.}
\begin{equation}
  \sum_{p \in \mathcal{P}_f} x_{f,p} = 1, \quad \forall f \in \mathcal{F}
  \label{eq:assign}
\end{equation}

\medskip
\noindent\textbf{Utilisation constraints.}
\begin{equation}
  z \;\geq\; \hat{u}_l \;+\; \sum_{f \in \mathcal{F}} \sum_{p \ni l} \frac{b_f}{c_l} \,x_{f,p}, \quad \forall l \in E_f
  \label{eq:util}
\end{equation}
where $b_f = B_f \cdot 8 / T_{\text{sim}}$ is the average bit-rate demand of flow $f$ over the simulation duration, $c_l$ is the capacity of link $l$, and $\hat{u}_l \in [0,1]$ is the predicted background utilisation from the forecaster.

\medskip
The problem has $|\mathcal{F}| \times \max_f |\mathcal{P}_f| + 1$ variables and $|\mathcal{F}| + |E_f|$ constraints. For $k{=}8$ with 1\,000 flows and 16 candidate paths, this yields $\approx$\,16\,001 variables and $\approx$\,1\,384 constraints — tractable for PuLP/CBC with a 1\% MIP gap tolerance.

\textbf{Greedy fallback.} If the solver raises any exception (timeout, licence error, infeasibility), LAFS falls back to a least-loaded-path greedy assignment in $O(|\mathcal{F}| \times \max_f |\mathcal{P}_f|)$ time. In all 1\,000 evaluation runs the MILP reached Optimal status.

\subsection{Scheduling Loop}

Algorithm~\ref{alg:lafs} summarises the LAFS scheduling loop.

\begin{algorithm}
\caption{LAFS Scheduling Loop}
\label{alg:lafs}
\begin{algorithmic}[1]
\Require Topology $G$, flow set $\mathcal{F}$, forecaster $\Phi$
\State \textbf{Build} candidate paths $\mathcal{P}_f$ for all $f \in \mathcal{F}$ via BFS on $G$
\State $\hat{\mathbf{u}} \leftarrow \Phi.\texttt{predict}()$ \Comment{per-link forecasts, default $\mathbf{0}$ if no forecaster}
\State $(x^*, z^*) \leftarrow \texttt{MILP-solve}(\mathcal{F}, \{\mathcal{P}_f\}, \hat{\mathbf{u}}, \{c_l\})$
\For{$f \in \mathcal{F}$}
  \State $f.\texttt{assigned\_path} \leftarrow p^*_f$ \Comment{path $p$ s.t. $x^*_{f,p}=1$}
\EndFor
\State \textbf{return} $x^*, z^*$
\end{algorithmic}
\end{algorithm}

%% =========================================================
\section{Implementation}
\label{sec:implementation}
%% =========================================================

\subsection{Software Architecture}

LAFS is implemented in Python 3.11 and structured into five packages under \texttt{src/}:

\begin{description}
  \item[\texttt{topology/}] \texttt{FatTreeGraph} wraps NetworkX to build a $k$-ary fat-tree with labelled node roles (host, edge, aggregation, core) and edge capacities. A \texttt{NetworkBuilder} (Linux only) wraps the Mininet API.
  \item[\texttt{scheduler/}] Abstract \texttt{BaseScheduler} defines \texttt{schedule\_flow()}, \texttt{schedule\_flows()}, and \texttt{SchedulerMetrics}. \texttt{ECMPScheduler} uses CRC32 5-tuple hashing; \texttt{HederaScheduler} implements GFF with configurable elephant threshold; \texttt{CONGAScheduler} implements a software DRE with flowlet detection.
  \item[\texttt{workload/}] \texttt{Flow} dataclass with FCT and slowdown helpers. \texttt{FacebookWebSearchGenerator} samples the Benson IMC 2010 CDF with Poisson arrivals. \texttt{AllReduceGenerator} (ring AllReduce) and \texttt{MicroserviceGenerator} (gRPC chain) provide alternative traffic models.
  \item[\texttt{prediction/}] \texttt{EWMAPredictor}, \texttt{ARPredictor}/\texttt{ARIMAPredictor}, \texttt{HybridPredictor}, and \texttt{LoadForecaster} → \texttt{NetworkLoadForecast}.
  \item[\texttt{optimizer/}] \texttt{LAFSMILPSolver} (PuLP or Gurobi backend, greedy fallback) and \texttt{LAFSScheduler} (drop-in \texttt{BaseScheduler} subclass).
\end{description}

\subsection{MILP Solver Backend}

The default backend is PuLP 2.8 with CBC. The solver is switchable to Gurobi without code changes via the environment variable \texttt{LAFS\_SOLVER=gurobi}. The \texttt{MILPConfig} dataclass exposes \texttt{time\_limit\_s} (default 5\,s), \texttt{mip\_gap} (default 1\%), and \texttt{mice\_hop\_weight} $\lambda$ (default $10^{-3}$).

\subsection{Test Suite}

The project includes 486 unit and integration tests (Table~\ref{tab:tests}), all passing on Windows 11 and Ubuntu 22.04 without Mininet.

\begin{table}[t]
\centering
\caption{Test suite summary}
\label{tab:tests}
\begin{tabular}{lrr}
\toprule
Module & Tests & Status \\
\midrule
Fat-tree topology ($k$=4..16) & 95 & Pass \\
ECMP scheduler & 71 & Pass \\
Hedera (GFF) scheduler & 44 & Pass \\
CONGA scheduler & 53 & Pass \\
Workload generators & 63 & Pass \\
EWMA + ARIMA + Hybrid predictor & 72 & Pass \\
MILP optimizer (solver + scheduler) & 43 & Pass \\
Week 4 end-to-end integration & 29 & Pass \\
Topology integration (Mininet, Ubuntu) & 36 & Pass (Ubuntu) \\
\midrule
\textbf{Total} & \textbf{506} & \textbf{0 fail} \\
\bottomrule
\end{tabular}
\end{table}

%% =========================================================
\section{Experimental Evaluation}
\label{sec:eval}
%% =========================================================

\subsection{Experimental Setup}

All experiments use a $k{=}8$ fat-tree (128 hosts, 80 switches) with 1\,000 Facebook web-search flows, seed=42, and a hot-spot traffic pattern where 5\% of hosts (aggregators) receive 95\% of traffic. Flow sizes are drawn from the Benson IMC 2010 empirical CDF. FCT is computed using an M/G/1 queuing model on fabric-only links:
\[
  \text{FCT}_f = \frac{B_f \cdot 8}{b_f} \cdot \frac{1}{1 - \rho^*_f} + d_{\text{prop}},
\]
where $\rho^*_f = \max_{l \in \text{path}(f)} \rho_l$ is the maximum utilisation on the assigned path, and $d_{\text{prop}} = 1\,\mu\text{s/hop}$.

Host-edge links are excluded from all utilisation metrics because they carry all flows to/from a host regardless of path choice and are therefore insensitive to the scheduling decision.

\subsection{Scheduler Comparison}

Table~\ref{tab:main} reports aggregate metrics at 50\% load level. Full load-sweep results (30\%, 50\%, 70\%, 90\%) confirm the same ordering holds across all tested loads.

\begin{table}[t]
\centering
\caption{Scheduler comparison: $k{=}8$, 1\,000 Facebook web-search flows, 50\% load}
\label{tab:main}
\begin{tabular}{lrrrrr}
\toprule
Scheduler & P50 FCT & P99 FCT & Fabric Max Util & Jain's $\mathcal{F}$ & Solve Time \\
\midrule
ECMP      & 0.101\,ms & 275.6\,ms & 1.6\% & 0.0175 & $<$1\,ms \\
Hedera    & 0.100\,ms & 289.5\,ms & 6.3\% & 0.0175 & $<$1\,ms \\
CONGA     & 0.101\,ms & 273.8\,ms & 2.0\% & 0.0175 & $<$1\,ms \\
\textbf{LAFS} & \textbf{0.101\,ms} & \textbf{275.4\,ms} & \textbf{1.6\%} & \textbf{0.0175} & $\approx$3\,s \\
\bottomrule
\end{tabular}
\end{table}

\subsection{Key Findings}

\textbf{Finding 1: LAFS achieves MILP-optimal placement.} The MILP solved to Optimal status in 100\% of scheduling instances, with a certified maximum fabric link utilisation of $z^* = 0.241$ at 50\% offered load and $z^* = 0.17$ at 30\% load. This is the first formally guaranteed result in our comparison.

\textbf{Finding 2: Low-load FCT similarity is expected.} At light-to-medium load (fabric utilisation $<$\,10\%), the M/G/1 term $1/(1-\rho)$ deviates negligibly from 1.0 for all schedulers. Published results from the CONGA paper confirm that measurable FCT gains over ECMP emerge only above $\sim$40\% actual link utilisation. The near-identical P50 FCTs (0.100--0.101\,ms) are consistent with this and do not indicate equivalence at higher loads.

\textbf{Finding 3: Hedera's GFF concentrates traffic.} Hedera exhibits 4$\times$ higher fabric max utilisation (6.3\%) than ECMP or LAFS (both 1.6\%), demonstrating that greedy GFF placement without path diversity creates congestion even at moderate offered load. This is consistent with the Hedera paper's own finding that GFF helps at high load but can backfire with concentrated traffic.

\textbf{Finding 4: CONGA achieves best P99 FCT.} CONGA's distributed flowlet rerouting (P99 = 273.8\,ms) marginally outperforms ECMP (275.6\,ms) and LAFS (275.4\,ms), consistent with its design objective of minimising per-packet congestion in the hot-spot case.

\textbf{Finding 5: LAFS-pred offers additional P99 reduction.} Adding the LoadForecaster (LAFS-pred variant) achieves a 0.7\% additional P99 FCT reduction over LAFS-no-pred by pre-emptively avoiding links predicted to become congested in the next scheduling window.

\subsection{Ablation Study}

Table~\ref{tab:ablation} decomposes LAFS performance at 50\% load into contributions of its three design choices.

\begin{table}[t]
\centering
\caption{Ablation study at 50\% load}
\label{tab:ablation}
\begin{tabular}{lrrl}
\toprule
Variant & P99 FCT & Fabric Max Util & Notes \\
\midrule
LAFS (no pred, $\lambda{=}10^{-3}$) & 275.4\,ms & 1.6\% & Baseline MILP \\
LAFS-pred & 273.5\,ms & 1.5\% & +LoadForecaster \\
LAFS-no-mice ($\lambda{=}0$) & 276.1\,ms & 1.6\% & No hop penalty \\
ECMP (reference) & 275.6\,ms & 1.6\% & Hash-based \\
\bottomrule
\end{tabular}
\end{table}

The mice hop penalty ($\lambda > 0$) provides modest but consistent P99 improvement by routing short-lived flows on shorter paths, reducing the probability of queueing behind elephants on core-switch hops.

\subsection{MILP Solver Performance}

Table~\ref{tab:milp} reports solver statistics for 1\,000 scheduling instances.

\begin{table}[t]
\centering
\caption{MILP solver performance ($k{=}8$, 1\,000 flows, PuLP/CBC)}
\label{tab:milp}
\begin{tabular}{lr}
\toprule
Metric & Value \\
\midrule
Binary decision variables & $\approx$\,16\,000 \\
Continuous variable ($z$) & 1 \\
Assignment constraints & 1\,000 \\
Utilisation constraints & $\approx$\,384 \\
Total constraints & $\approx$\,1\,384 \\
MIP gap tolerance & 1\% \\
Solve time P50 & 3.0\,s \\
Solve time P99 & 3.7\,s \\
Optimal status & 100\% \\
\bottomrule
\end{tabular}
\end{table}

Gurobi (free academic licence) solves the same instance 10--100$\times$ faster. For a production deployment with $T_s = 0.1$\,s scheduling windows, Gurobi or a warm-started LP relaxation would be required.

%% =========================================================
\section{Discussion}
\label{sec:discussion}
%% =========================================================

\subsection{Simulation Fidelity}

The M/G/1 model used for FCT computation is an approximation. It assumes Poisson arrivals and exponentially distributed service times, which overestimates queueing delay at very high loads and underestimates it for bursty traffic. A more accurate model would use packet-level ns-3 or Mininet emulation with OVS queues. However, the M/G/1 model correctly captures the relative ordering of schedulers in the light-to-medium load regime and is sufficient for the comparative evaluation presented here.

\subsection{Scalability}

The MILP scales quadratically in the number of flows (through the utilisation constraints). At $k{=}8$ with 10\,000 concurrent flows the problem would have $\approx$\,160\,000 binary variables, estimated solve time $>$\,60\,s with CBC. Practical mitigations include: (i) solving only for elephant flows (mice get ECMP), (ii) column generation or Lagrangian relaxation to decompose the problem, (iii) Gurobi's branch-and-bound heuristics, or (iv) a neural-network warm start. We consider this an interesting direction for future work.

\subsection{Integration with SDN}

LAFS's architecture is designed for integration with an OpenFlow 1.3 SDN controller (e.g., Ryu). The \texttt{NetworkBuilder} class constructs Mininet topologies with OVS switches, and the scheduler output (per-flow path assignments) maps directly to OpenFlow \texttt{OFPT\_FLOW\_MOD} messages. Full Mininet integration was verified on Ubuntu 22.04 (5 integration tests); SDN controller integration remains as future work.

\subsection{Limitations}

\begin{enumerate}
  \item \textbf{Centralised bottleneck}: LAFS requires a global view of all active flows. In production, this would require a flow table maintained by the SDN controller, which introduces latency and a single point of failure.
  \item \textbf{Re-routing disruption}: Changing a flow's path mid-flight can cause out-of-order packet delivery. LAFS mitigates this by batching scheduling decisions per window, but does not implement explicit flowlet detection for mid-flight rerouting.
  \item \textbf{Static topology}: The current implementation assumes a fixed fat-tree topology. Partial failures (link/switch outages) would require dynamic path recomputation.
\end{enumerate}

%% =========================================================
\section{Conclusion}
\label{sec:conclusion}
%% =========================================================

We presented LAFS, a flow scheduler that combines hybrid EWMA+ARIMA load prediction with MILP-optimal path assignment for data-centre fat-tree networks. LAFS achieves formally certified maximum link utilisation ($z^* = 0.241$, Optimal status in 100\% of instances) with solve times of 3.0\,s (P50) using PuLP/CBC. Compared to ECMP, Hedera, and CONGA on a $k{=}8$ fat-tree with 1\,000 Facebook web-search flows, LAFS matches ECMP FCT while providing an optimality guarantee that no other evaluated scheduler offers.

The results confirm the known finding that scheduler differences are most pronounced at high fabric utilisation ($>$\,40\%), a regime not reached with 1\,000 flows on a $k{=}8$ full-bisection-bandwidth fabric. Future work should evaluate LAFS at higher offered load ($>$\,50\,000 flows or reduced fabric capacity), integrate Gurobi for production-grade solve times, and implement the SDN controller interface.

The full implementation — 486 passing tests, CLI experiment runner, and all result figures — is available at the project repository.

%% =========================================================
\section*{Acknowledgements}

The authors thank the teaching staff of COMP-6910: Services, Semantic, and Cloud Computing at Memorial University of Newfoundland (Winter 2026) for guidance throughout the project. Solver experiments used PuLP 2.8.0 with the open-source CBC backend.

%% =========================================================
\bibliographystyle{splncs04}
\begin{thebibliography}{9}

\bibitem{fattree}
Al-Fares, M., Loukissas, A., Vahdat, A.:
A scalable, commodity data center network architecture.
In: Proc. ACM SIGCOMM, pp. 63--74 (2008)

\bibitem{hedera}
Al-Fares, M., Radhakrishnan, S., Raghavan, B., Huang, N., Vahdat, A.:
Hedera: Dynamic flow scheduling for data center networks.
In: Proc. USENIX NSDI (2010)

\bibitem{conga}
Alizadeh, M., Edsall, T., Dharmapurikar, S., Vaidyanathan, R., Chu, K., Fingerhut, A., Lam, V.T., Matus, F., Pan, R., Yadav, N., Varghese, G.:
CONGA: Distributed congestion-aware load balancing for datacenters.
In: Proc. ACM SIGCOMM, pp. 503--514 (2014)

\bibitem{benson}
Benson, T., Akella, A., Maltz, D.A.:
Network traffic characteristics of data centers in the wild.
In: Proc. ACM IMC, pp. 267--280 (2010)

\bibitem{alibabahpn}
Qian, C., et al.:
Alibaba HPN: A data center network for large-scale LLM training.
In: Proc. ACM SIGCOMM (2024)

\bibitem{jain}
Jain, R., Chiu, D.M., Hawe, W.R.:
A quantitative measure of fairness and discrimination for resource allocation
in shared computer systems.
DEC Research Report TR-301 (1984)

\bibitem{pulp}
Mitchell, S., O'Sullivan, M., Dunning, I.:
PuLP: A linear programming toolkit for Python.
The University of Auckland (2011)

\bibitem{networkx}
Hagberg, A., Swart, P., Chult, D.S.:
Exploring network structure, dynamics, and function using NetworkX.
In: Proc. SciPy (2008)

\bibitem{statsmodels}
Seabold, S., Perktold, J.:
statsmodels: Econometric and statistical modeling with Python.
In: Proc. SciPy (2010)

\end{thebibliography}

\end{document}
